% ==============================================================
% Kapitel 4: Ziele
% Autoren: beide zusammen
% Umfang: 1 Seite
% ==============================================================

\section*{Ziele}%
\label{sec:zielsetzung}

Das Ziel dieser Diplomarbeit ist es, ein wartbares, erweiterbares und einfach 
aufsetzbares Zeitmesssystem für die \textit{Zero Emission Challenge} zu entwickeln, 
welches den gesamten Ablauf von der Zeiterfassung über die Punkteauswertung bis zur Ergebnisanzeige umfasst.
Das System soll zusätzlich auch die Verwaltung von Teams und Benutzern ermöglichen und 
eine rollenbasierte Zugriffskontrolle implementieren.

Das entwickelte System soll:
\begin{itemize}
    \item Renndaten aller Bewerbe zuverlässig erfassen, speichern und auswerten
    \item Zeitstempel von Lichtschranken per MQTT empfangen und in der 
          Zeitnehmer-Applikation zur Kontrolle und manuellen Korrektur bereitstellen
    \item Eine webbasierte Plattform zur Verwaltung von Teams, Benutzern und 
          Ergebnissen bereitstellen
    \item Sicherheit durch rollenbasierte Zugriffskontrolle und eine zentrale 
          Authentifizierung über Keycloak gewährleisten
    \item Durch eine Microservice-Architektur eine hohe Verfügbarkeit und 
          horizontale Skalierbarkeit ermöglichen
    \item Mittels Docker und Kubernetes einfach und plattformunabhängig 
          betreibbar sein
    \item Softwarequalität durch automatisierte Tests und CI/CD-Pipelines 
          sicherstellen
\end{itemize}